\subsection{LINEAR FORMS; DUAL OF A MODULE}

Let \(\mathbf{E}\) be a left \(\mathbf{A}\) -module. As \(\mathbf{A}\) is an \((\mathbf{A}, \mathbf{A})\) -bimodule, \(\operatorname{Hom}_{\mathbf{A}}(\mathbf{E}, \mathbf{A}_s)\) has a canonical right \(\mathbf{A}\) -module structure (§ 1, no. 14).

DEFINITION 2. For every right A-module E, the right A-module \(\mathrm{Hom}_{\mathbf{A}}(\mathbf{E},\mathbf{A}_{s})\) is called the dual module of E (or simply the dual† of E) and its elements are called the linear forms on E.

If E is a right A-module, the set \(\operatorname{Hom}_{\mathbf{A}}(\mathbf{E}, \mathbf{A}_{d})\) with its canonical left A-module structure is likewise called the dual of E and its elements are called linear forms on E.

In this chapter, E* will be used to denote the dual of a (left or right) A-module E.

Example. *On the vector space (with respect to the field \(\mathbf{R}\) ) of continuous real-valued functions on an interval \([a, b]\) of \(\mathbf{R}\) , the mapping \(x \mapsto \int_{a}^{b} x(t) dt\) is a linear form.*

Let E be a left A-module and \(E^{*}\) its dual; for every ordered pair of elements \(x \in E\) , \(x^{*} \in E^{*}\) , the element \(x^{*}(x)\) of A is denoted by \(\langle x, x^{*} \rangle\) . Then the relations

\begin{enumerate}
\def\labelenumi{(\arabic{enumi})}
\setcounter{enumi}{8}
\tightlist
\item
\end{enumerate}

\[
\langle x + y, x ^ {*} \rangle = \langle x, x ^ {*} \rangle + \langle y, x ^ {*} \rangle\tag{10}
\]

\[
\langle x, x ^ {*} + y ^ {*} \rangle = \langle x, x ^ {*} \rangle + \langle x, y ^ {*} \rangle\tag{11}
\]

\[
\langle \alpha x, x ^ {*} \rangle = \alpha \langle x, x ^ {*} \rangle\tag{12}
\]

\[
\langle x, x ^ {*} \alpha \rangle = \langle x, x ^ {*} \rangle \alpha
\]

hold for \(x, y\) in \(\mathbf{E}, x^*, y^*\) in \(\mathbf{E}^*\) and \(\alpha \in \mathbf{A}\) . The mapping \((x, x^*) \mapsto \langle x, x^* \rangle\) of \(\mathbf{E} \times \mathbf{E}^*\) into \(\mathbf{A}\) is called the canonical bilinear form on \(\mathbf{E} \times \mathbf{E}^*\) (the notion of bilinear form will be defined generally in IX, § 1). Every linear form \(x^*\) on \(\mathbf{E}\) can be considered as the partial mapping \(x \mapsto \langle x, x^* \rangle\) corresponding to the canonical bilinear form.

When E is a right A-module, the value \(x^{*}(x)\) of a linear form \(x^{*} \in E^{*}\) at an element \(x \in E\) is denoted by \(\langle x^{*}, x \rangle\) and the formulae corresponding to (11) and (12) are written as

\[
\begin{array}{l} \langle x ^ {*}, x \alpha \rangle = \langle x ^ {*}, x \rangle \alpha \\ \langle \alpha x ^ {*}, x \rangle = \alpha \langle x ^ {*}, x \rangle . \end{array}
\]

When A is commutative, either notation is permissible.

PROPOSITION 5. For every ring A, the mapping which associates with every \(\xi\in\mathbf{A}\) the linear form \(\eta\mapsto\eta\xi\) on \(\mathbf{A}_{s}\) is an isomorphism of \(\mathbf{A}_{d}\) onto the dual of \(\mathbf{A}_{s}\) .

It is the particular case of the canonical isomorphism \(E \to \operatorname{Hom}_{A}(A_{s}, E)\) of §1, no. 14, Remark 2, corresponding to \(E = A_{s}\) ; the commutativity of diagram (50) of §1, no. 14, shows that we have here an isomorphism of right A-modules.

If \(A_{d}\) is identified with the dual of \(A_{s}\) by means of the isomorphism of Proposition 5, the canonical bilinear form on \(A_{s} \times A_{d}\) is then expressed by

\[
\langle \xi , \xi^ {*} \rangle = \xi \xi^ {*} \quad \text { for } \xi , \xi^ {*} \text { in } A.\tag{13}
\]

Similarly, the dual of \(A_{d}\) is canonically identified with \(A_{s}\) , the canonical bilinear form on \(A_{d} \times A_{s}\) being expressed by

\[
\langle \xi^ {*}, \xi \rangle = \xi^ {*} \xi \quad \text { for } \xi , \xi^ {*} \text { in } A.\tag{14}
\]

